\section{Observation, Memory, Control, and Time in Finite Models}
\label{op:chapter}

This section studies what a specified collection of observations determines about
a finite mathematical model.  Its objects are conventional probability kernels,
finite hidden-state models, quantum states, and continuous-time Markov generators.
An observation law, an intervention, and a calibration premise are distinct
inputs.  Equality of observable laws will always refer to an explicitly stated
experimental menu.  Identifiability within such a menu does not establish that
the chosen states, controls, environmental factors, or time coordinate describe
a particular physical system.

We use row probability vectors.  For probability measures on a finite set,
\(\operatorname{TV}(p,q)=\frac12\sum_x|p_x-q_x|\).
Matrix \(1\)-norms below are entrywise unless specified otherwise; the Frobenius
norm is denoted by \(\|\cdot\|_F\).  The vector \(\mathbf1\) has all entries one.
All logarithms in likelihood calculations are natural; binary entropy is
explicitly measured in bits.  Population probabilities are exact mathematical
inputs unless a finite-sample statement is given.

\subsection{Controlled reduction and the limits of short records}
\label{op:reduction}

Let \(X,Z\) be finite sets, \(q:X\to Z\) a surjection, and \(P_a\) a
row-stochastic transition matrix on \(X\) for each allowed action \(a\).
The full process is Markov conditional on the selected action.  A policy may
randomize and may depend on the complete visible history, but it has no extra
model-specific hidden-state information.  Write
\[
 M_a(x,z')=\sum_{y:q(y)=z'}P_a(x,y).
\]
The familiar strong-lumpability condition is the action-free version of the
following statement; see Wolfer and Watanabe~\cite{op:wolfer}.

\begin{theorem}[Exact controlled reduction]\label{op:lumpability}
There are kernels \(Q_a\) on \(Z\) that reproduce the complete visible-history
law for every initial distribution on \(X\) and every common visible-history
policy if and only if, for every action \(a\), \(M_a(x,\cdot)\) is constant
on each fiber of \(q\).  In that event the common row is
\(Q_a(q(x),\cdot)\).
\end{theorem}
\begin{proof}
For necessity, choose point initial distributions at two states \(x,x'\) in
one fiber and apply the same fixed action.  Both reduced initial states are
\(q(x)\); their next-observation laws must therefore be equal.
For sufficiency, condition on any visible history of positive probability
and the next selected action.  The conditional law of the hidden state is
supported in the fiber of the current observation.  Mixing identical pushed
rows over that posterior gives the stated \(Q_a\) row.  The policy uses the
same conditional distribution of actions in both descriptions.  Induction
over the finite history proves equality, including the recorded actions.
\end{proof}

If only specified starting states are admitted, the fiber condition need
only be imposed on the appropriately reachable domain.  Reduction for one
special initial law is a weaker property and does not imply the all-start
statement.

\begin{example}[Fair pairs with predictive memory]\label{op:four-state}
On states \(00,01,10,11\), observe the final bit and use
\[
 P=\begin{pmatrix}
 3/4&1/4&0&0\\
 0&0&1/4&3/4\\
 1/4&3/4&0&0\\
 0&0&3/4&1/4
 \end{pmatrix}.
\]
The uniform distribution is stationary because each column sums to one.
Both append choices have positive probability; the shift graph is strongly
connected and has self-loops, so the chain is irreducible and aperiodic.
In stationarity every adjacent observed pair has probability \(1/4\).
Nevertheless,
\[
 P(Y_1=a,Y_2=b,Y_3=c)
 =\frac18\left[1+\frac12(-1)^{a+b+c}\right],
 \qquad
 \operatorname{TV}(P_{Y_{1:3}},\operatorname{Unif}\{0,1\}^3)=\frac14 .
\]
For example, the next-zero probabilities after \(00\) and \(10\) are \(3/4\)
and \(1/4\), respectively.  The visible process is not first-order Markov.
The best uniform row approximation by a kernel depending on the last bit
has error \(1/4\): the two relevant Bernoulli rows have distance \(1/2\),
and their fair midpoint attains half that distance.
\end{example}

Hidden state is not by itself an obstruction.  For example,
\[
 P((z,h),(z',h'))=Q(z,z')\,\mathbf1_{\{h'=1-h\}}
\]
has a perpetually alternating hidden bit and an exactly autonomous visible
kernel \(Q\).

\begin{theorem}[Uniform row error controls complete histories]
\label{op:history-bound}
Suppose, uniformly over all admitted \(x,a\),
\[
 \operatorname{TV}\bigl(M_a(x,\cdot),Q_a(q(x),\cdot)\bigr)\le\delta,
 \qquad 0\le\delta\le1.
\]
Start with the same visible law and use the same visible-history policy.
After \(N\) transitions the total variation distance between complete
visible histories, including actions, is at most
\[
 1-(1-\delta)^N\le N\delta .
\]
The first bound is sharp.
\end{theorem}
\begin{proof}
Couple the initial observations identically.  As long as the histories
agree, use the same policy randomness.  The next visible law of the full
model is a posterior mixture of pushed rows, hence is within \(\delta\)
of the reduced row by convexity of total variation.  A maximal coupling
keeps the next observations equal with probability at least \(1-\delta\).
Consequently the coupled histories agree with probability at least
\((1-\delta)^N\).  The coupling inequality proves the claim.
For sharpness, let the reference chain remain at state \(0\), with state
\(1\) absorbing if ever entered.  Let the comparison chain escape from
\(0\) to absorbing \(1\) with probability \(\delta\) on each step.
The row premise holds at both states.  The reference history is all zero,
and its probability under the comparison model is exactly
\((1-\delta)^N\).
\end{proof}

Thus a test of size \(\alpha\) for the reduced history law has size at most
\(\min\{1,\alpha+1-(1-\delta)^N\}\) for the full model.  An average
stationary row error supplies no such uniform conditional premise.

An intervention on fine states need not define an intervention on coarse
states.  As a simple exact example, let \(q(1)=q(2)=0\), \(q(3)=q(4)=1\).
The point masses at \(1,2\) have identical coarse images.  A permutation
interchanging \(2,3\) while fixing \(1,4\) sends them to distinct coarse
point masses.  No map of the original coarse distribution can represent
that operation on both inputs.  The quantum counterpart is equally direct:
\(Z\)-basis dephasing sends both \(|+\rangle\langle+|\) and
\(|-\rangle\langle-|\) to \(I/2\), while a Hadamard operation followed by
that dephasing distinguishes them perfectly.  Specifying the admitted
interventions is therefore part of specifying a reduced model.

\begin{example}[A full-support non-descent witness]\label{op:non-descent}
For the same partition \(\{1,2\},\{3,4\}\), take the strictly positive
fine laws
\[
 p=(3/5,3/20,1/5,1/20),\qquad
 p'=(3/20,3/5,1/5,1/20).
\]
Both have coarse law \((3/4,1/4)\).
Interchanging the second and third fine labels sends them to coarse
laws \((4/5,1/5)\) and \((7/20,13/20)\), of total variation \(9/20\).
Thus the failure is not confined to boundary point masses.
Exactly eight of the twenty-four fine permutations descend for all
fine distributions: the permutation must send each whole input block
to one output block, giving two block permutations and two internal
permutations in each block, hence \(2(2!)^2=8\).
\end{example}

\subsection{A parity-shift process and its exact innovation coordinates}
\label{op:parity}

Fix \(L\ge1\) and \(-1<\eta<1\).  The hidden state is the last \(L\) bits,
initially uniform.  From \(x=(x_1,\ldots,x_L)\), append \(b\) with
probability
\begin{equation}\label{op:parity-kernel}
 P_\eta\bigl(x,(x_2,\ldots,x_L,b)\bigr)
 =\frac{1+\eta(-1)^{x_1+\cdots+x_L+b}}2 .
\end{equation}
This family is a parity form of the two-faced chains studied by
Ryabko~\cite{op:ryabko}, with parameter \(p=(1+\eta)/2\).

\begin{proposition}[Finite-depth agreement]\label{op:finite-depth}
The uniform law on \(L\)-bit states is stationary.  Every contiguous
observed block of length at most \(L\) is uniform, whereas a block of
length \(L+1\) has law
\[
 2^{-(L+1)}\left[1+\eta(-1)^{\sum_{j=1}^{L+1}y_j}\right].
\]
Its distance from the uniform law is \(|\eta|/2\).
For \(\eta\ne0\), the observed process has Markov order exactly \(L\).
\end{proposition}
\begin{proof}
Each destination state has two predecessors, differing in their dropped
bit.  Their transition probabilities sum to one.  Thus a uniform state
distribution remains uniform.  Marginalization gives every shorter
contiguous block law.  Multiplying the initial \(2^{-L}\) probability by
the append probability gives the displayed \(L+1\) law, whose \(2^{L+1}\)
absolute discrepancies are all \(2^{-(L+1)}|\eta|\).
The transition rule depends on no more than the last \(L\) bits.
For \(\eta\ne0\), changing the oldest of those bits while keeping the last
\(L-1\) fixed changes the next-bit probability.  All such histories have
positive probability, proving that smaller order is impossible.
\end{proof}

The statement concerns contiguous blocks.  It does not assert independence
of every collection of \(L\) widely separated observations.  Since \(L\)
can be arbitrarily large, no common finite validation depth distinguishes
the whole family from fair independent bits.

For a record of length \(L+m\), define parity innovations
\[
 E_t=Y_t\mathbin{\oplus}\cdots\mathbin{\oplus}Y_{t+L},
 \quad 1\le t\le m,
\]
where \(\oplus\) is addition modulo two.

\begin{theorem}[Innovation factorization]\label{op:innovations}
The map from \(Y_{1:L+m}\) to \((Y_{1:L},E_{1:m})\) is bijective and
\[
 P(Y_{1:L}=x,E_{1:m}=e)
 =2^{-L}\prod_{t=1}^m\frac{1+\eta(-1)^{e_t}}2.
\]
Consequently the innovations are independent and identically distributed,
with \(P(E_t=0)=(1+\eta)/2\), and are independent of the initial block.
\end{theorem}
\begin{proof}
The recursion
\(Y_{t+L}=E_t\oplus Y_t\oplus\cdots\oplus Y_{t+L-1}\)
recovers the record uniquely.  Its probability is the initial block
probability times the \(m\) append probabilities in
\eqref{op:parity-kernel}.  Each append factor depends only on \(e_t\),
giving the asserted product law.  Bijection alone would not establish
independence; the factorization of probabilities is essential.
\end{proof}

Writing \(h_2\) for binary entropy in bits and \(d_{\rm Ber}\) for
Bernoulli relative entropy in nats, this gives
\[
 H(Y_{1:L+m})=L+m h_2((1+\eta)/2),\qquad
 D(P_{Y_{1:L+m}}\Vert\operatorname{Unif})=
 m d_{\rm Ber}((1+\eta)/2\Vert1/2).
\]
These are entropy and divergence per indexed observation, not a physical
communication capacity or a rate per unit of laboratory time.
Cancellation of overlapping parities also gives
\[
 (-1)^{Y_t+Y_{t+L+1}}=(-1)^{E_t+E_{t+1}},\qquad
 P(Y_t=a,Y_{t+L+1}=b)=\tfrac14[1+\eta^2(-1)^{a+b}].
\]
The two retained outcomes therefore have distance \(\eta^2/2\) from a
fair independent pair, although their acquisition spans \(L+2\) process
positions.  Repeated overlapping sparse pairs are not automatically
independent trials.
For \(L\ge2\), pairs at lags \(1,\ldots,L\) are independent and fair:
for lag \(L\), marginalizing any other coordinate of the \(L+1\) parity
law removes its perturbation, and smaller lags follow from the uniform
block laws.  The qualifier \(L\ge2\) is necessary.

\subsubsection{Ideal fixed-horizon testing}

\begin{lemma}[Finite likelihood-ratio ordering]\label{op:np}
For two strictly positive probability laws \(p_0,p_1\) on a finite
sample space, a randomized test of prescribed size \(\alpha\) has maximum
power by rejecting in decreasing order of \(p_1/p_0\), randomizing only
on a boundary level.
\end{lemma}
\begin{proof}
If the threshold is \(c\) and the proposed test is \(\phi_*\), then
\((\phi_*-\phi)(p_1-cp_0)\ge0\) pointwise for every other test
\(0\le\phi\le1\).  Summation and equal size imply
\(\sum\phi_*p_1\ge\sum\phi p_1\).  A boundary randomization fills the
required size.  This is the finite Neyman--Pearson argument
\cite{op:neyman}.
\end{proof}

With known \(L\), error-free readout, and alternatives \(\eta=0\) versus
\(\eta=1/2\), let \(C_m=\sum_{t=1}^m\mathbf1_{\{E_t=0\}}\).
The likelihood ratio of the complete record is \(3^{C_m}/2^m\);
the initial block is ancillary.  Thus
\[
 C_m\sim\operatorname{Bin}(m,1/2)\quad\hbox{or}\quad
 \operatorname{Bin}(m,3/4).
\]
At size \(1/20\), the optimum powers for \(m=39,40\) are respectively
\(0.9498706279714081\ldots\) and \(0.9530027582141789\ldots\).
At \(m=40\), reject for \(C_m\ge26\) and, when \(C_m=25\), reject with
probability \(4423136507/16760560440\).
The exact power is
\[
 \frac{720068525355674046204729}{755578637259143234191360}.
\]
Optimal full-record power cannot decrease with the horizon, because an
extra observation can be discarded.  The failure at \(39\) and success
at \(40\) therefore establish the minimum within this randomized
fixed-horizon experiment.  It uses \(L+40\) raw outputs.  The nonrandomized
rule \(m=42,C_m\ge27\) is a sufficient alternative, with size
\(0.04421477349\ldots\) and power \(0.95837130194\ldots\); no global
nonrandomized minimum is asserted.

\subsubsection{Independent readout errors create dependent innovations}

Now observe \(\widetilde Y_t=Y_t\oplus N_t\), where the \(N_t\) are iid
Bernoulli\((\epsilon)\), independent of the entire source, and
\(0\le\epsilon<1/2\).  Put \(a=1-2\epsilon\),
\(W_t=(-1)^{E_t}\), \(V_j=(-1)^{N_j}\), and
\[
 Z_t=(-1)^{\widetilde E_t}
     =W_t\prod_{j=t}^{t+L}V_j.
\]

\begin{proposition}[Noisy innovation law]\label{op:noisy-cov}
The mean and covariance function are
\[
 \mu=\eta a^{L+1},\quad \Gamma_0=1-\mu^2,\quad
 \Gamma_h=
 \begin{cases}
 \eta^2(a^{2h}-a^{2L+2}),&1\le h\le L,\\
 0,&h\ge L+1.
 \end{cases}
\]
Collections of \(Z\)'s separated by more than \(L\) indices are
independent.  For \(C_m=\sum_{t=1}^m(1+Z_t)/2\),
\[
 \operatorname{Var}(C_m)=\frac14
 \left[m\Gamma_0+
 2\sum_{h=1}^{\min\{L,m-1\}}(m-h)\Gamma_h\right].
\]
\end{proposition}
\begin{proof}
The \(W_t\) are iid of mean \(\eta\), while the independent \(V_j\) have
mean \(a\).  In \(Z_tZ_{t+h}\), the \(L+1-h\) common detector factors
square to one when \(h\le L\), leaving \(2h\) factors and expectation
\(\eta^2a^{2h}\).  For larger \(h\), the defining independent variables
are disjoint, giving expectation \(\mu^2\) and the stronger setwise
independence statement.  Summing the covariances gives the variance.
\end{proof}

The observed initial \(L\)-bit block remains uniform and independent of
the complete observed innovation sequence: condition on all innovations
and all detector bits, and use the uniform independent source initial
block.  Thus the ordered noisy innovations are sufficient within this
specified family.  Their count generally is not.  For
\(L=2,\eta=1/2,\epsilon=1/10\),
\[
 P(\widetilde E_{1:3}=001)=681/5000,\qquad
 P(\widetilde E_{1:3}=010)=609/5000,
\]
although both words have the same count.  At \(m=40\), the count mean is
\(628/25\) and its variance \(185736/15625=11.887104\), rather than the
attenuated-binomial value \(9.34464\).

\begin{proposition}[Exact noisy count recursion]\label{op:count-recursion}
Let \(A_t(h,k)\) be the joint probability of the most recent \(L\) detector
bits \(h\) and \(k\) even observed innovations after \(t\) steps.
Initialize
\[
 A_0(h,0)=\epsilon^{|h|}(1-\epsilon)^{L-|h|}.
\]
For \(b\in\{0,1\}\), put \(r_b=\epsilon^b(1-\epsilon)^{1-b}\),
\(h'=(h_2,\ldots,h_L,b)\), and
\(q(h,b)=[1+\eta(-1)^{|h|+b}]/2\).
For each current mass, add
\[
 A_t(h,k)r_bq(h,b)\ \hbox{to } A_{t+1}(h',k+1),\qquad
 A_t(h,k)r_b[1-q(h,b)]\ \hbox{to } A_{t+1}(h',k).
\]
Then \(P(C_m=k)=\sum_h A_m(h,k)\).
\end{proposition}
\begin{proof}
The next detector bit is independent of the current detector history.
The next source innovation is independent of all preceding innovations
and all detector bits.  Conditional on the detector-window parity, its
observed even probability is exactly \(q(h,b)\).
The displayed update is the law of total probability on these two
new independent inputs.  Induction proves the result.
\end{proof}

For rational parameters this is an integer recurrence after keeping one
common denominator at each step.  Applying Lemma~\ref{op:np} to the count
laws, sorting their actual likelihood ratios rather than presuming their
order, gives the following finite certificate for
\(L=2,\eta=1/2,\epsilon=1/10\) against the fair iid null.
Every \(m=1,\ldots,182\) has optimal count-only power below \(0.95\) at
size \(0.05\); \(m=183\) has power \(0.95047965645\ldots\).
At \(m=182\) it is \(0.94920331752\ldots\).
An easily specified sufficient rule at \(183\) rejects for \(C_m\ge104\)
and with probability \(43/50\) at \(C_m=103\); its size and power are
\(0.04986347138\ldots\) and \(0.95033569239\ldots\).
The exact recurrence and rational comparisons are supplied with the
computational materials.  This minimum is for tests of \(C_m\) only.
Unlike the complete record, \(C_{m+1}\) need not recover \(C_m\);
monotonicity of the count experiments is not used.
Here \(183\) innovations require \(185\) raw outputs.
A nonrandomized sufficient rule uses \(188\) innovations and rejects
for \(C_m\ge106\); its size is \(0.04659560696\ldots\) and power
\(0.95203135112\ldots\).  No nonrandomized minimum is claimed.

The ideal \(40\)-innovation rule has power only
\(0.49258836198\ldots\) under this noisy alternative, while retaining its
null size.  Replacing its count law by a binomial with the attenuated
mean is incorrect.  If instead parity windows start \(L+1\) indices
apart, their defining source innovations and detector windows are
disjoint, so the selected parities are iid.  That thinning consumes
additional raw outputs.

At the population level, known \(L\), \(\eta\ne0\), and \(a>0\) allow
recovery from \(\mu=E Z_t\) and \(\nu=E Z_tZ_{t+1}\):
\[
 a=(\mu^2/\nu)^{1/(2L)},\qquad
 \epsilon=(1-a)/2,\qquad \eta=\mu/a^{L+1}.
\]
Indeed \(\mu^2/\nu=a^{2L}\) and \(\nu=\eta^2a^2\).
The cases \(\eta=0\) and \(a=0\) are degenerate, and the inversion is
poorly conditioned near them.  These identities do not supply a
finite-data confidence interval or validate the iid detector model.

\subsection{Classical realization rank and finite-sample certificates}
\label{op:rank}

Let \(H_L(x,y)=P(Y_{1:L}=x,Y_{L+1:2L}=y)\).
Any \(D\)-state classical hidden realization with a sufficient state at
this boundary factors the matrix as
\[
 H_L(x,y)=\sum_{s=1}^D P(Y_{1:L}=x,S_L=s)
                    P(Y_{L+1:2L}=y\mid S_L=s).
\]
Hence \(\operatorname{rank}H_L\le D\).
Every source, detector, or controller variable that carries predictive
information across the cut belongs in that sufficient state.  General
hidden-model realization theory distinguishes this necessary rank
condition from the existence of a positive realization
\cite{op:ohta,op:vanluyten}.

\begin{theorem}[Parity-block spectrum and exact state count]
\label{op:parity-rank}
In the ideal parity-shift model the singular values of \(H_L\) are
\[
 2^{-L}|\eta|^j,\qquad j=0,\ldots,L,
\]
with multiplicity \(\binom Lj\).
For \(\eta\ne0\), the minimum number of states of an exact classical
transition-emitting realization is \(2^L\).  It is one when \(\eta=0\).
\end{theorem}
\begin{proof}
Define separate bijections of the row and column bit strings by
\(u_t=x_t\oplus\cdots\oplus x_L\) and
\(v_t=y_1\oplus\cdots\oplus y_t\).
The \(t\)th innovation following the cut is \(u_t\oplus v_t\).
The innovation factorization therefore gives, up to separate row and
column permutations,
\[
 H_L=2^{-L}B_\eta^{\otimes L},\qquad
 B_\eta=\frac12\begin{pmatrix}1+\eta&1-\eta\\1-\eta&1+\eta\end{pmatrix}.
\]
The singular values of \(B_\eta\) are \(1,|\eta|\); tensor products give
the asserted list.  Separate row and column permutations preserve
singular values, though they need not be a similarity transformation.
For nonzero \(\eta\) the rank is \(2^L\), providing the lower bound.
The shift-state construction \eqref{op:parity-kernel} has exactly that
many states and provides the matching positive upper realization.
For \(\eta=0\) the outputs are iid fair bits.
\end{proof}

This \(2^L\) result belongs to the parity-shift model.  A stationary
random-phase concatenation of perturbed \(k\)-bit blocks is a different
process, with a different \(2k-1\) realization result.  The parameters
\(L,k\), their record laws, and their state counts are not interchangeable.
Neither count is a bound on physical constituents or on quantum memory.

\begin{theorem}[Readout and the noisy Walsh spectrum]\label{op:noisy-rank}
Independent symmetric readout with \(a=1-2\epsilon>0\) preserves the
rank \(2^L\) for \(\eta\ne0\).
For \(s\in\{0,1\}^L\), define
\[
 r_i(s)=\bigoplus_{j=1}^i s_j,\qquad
 \ell_i(s)=\bigoplus_{j=i}^L s_j .
 \]
The noisy matrix has singular values, indexed by all \(s\),
\[
 \sigma_s=2^{-L}|\eta|^{|s|}
                  a^{|r(s)|+|\ell(s)|}.
\]
In particular \(\sigma_{\min}\ge (|\eta|a^2/2)^L>0\).
\end{theorem}
\begin{proof}
Let \(D_\epsilon=\left(\begin{smallmatrix}1-\epsilon&\epsilon\\
\epsilon&1-\epsilon\end{smallmatrix}\right)\).
Independent noise on both blocks gives
\(\widetilde H_L=D_\epsilon^{\otimes L}H_L
(D_\epsilon^{\otimes L})^T\).
Its two eigenvalues are \(1,a\), so it is invertible for \(a>0\).
For the full spectrum, expand each transition factor to obtain
\[
 H_L(x,y)=2^{-2L}\sum_s\eta^{|s|}
             (-1)^{r(s)\cdot x+\ell(s)\cdot y}.
\]
The maps \(s\mapsto r(s)\), \(s\mapsto\ell(s)\) are bijections.
The normalized Walsh vectors
\(2^{-L/2}(-1)^{v\cdot x}\) form orthonormal bases.
In these bases, \(H_L\) is a weighted permutation matrix with weights
\(2^{-L}\eta^{|s|}\).  Readout multiplies its row and column Walsh
coefficients by \(a^{|r(s)|}\) and \(a^{|\ell(s)|}\).
The absolute weights are the singular values.  Since
\(|s|\le L\) and \(|r(s)|+|\ell(s)|\le2L\), the lower bound follows.
\end{proof}

The same \(2^L\)-state realization with memoryless noisy emission gives
the upper bound after readout.  At \(\epsilon=1/2\) the observed process
is iid fair and the rank collapses.  Exact invertibility does not imply
a well-conditioned statistical inversion.

\begin{lemma}[Singular-tail distance]\label{op:tail}
For a real matrix \(H\) with singular values in decreasing order, define
\(T_D(H)=\sum_{j>D}\sigma_j(H)^2\).
Then
\[
 d_D(H):=\inf_{\operatorname{rank}Q\le D}\|H-Q\|_F
       =\sqrt{T_D(H)},\qquad
 |d_D(H)-d_D(G)|\le\|H-G\|_F.
\]
\end{lemma}
\begin{proof}
If \(Q\) has rank at most \(D\), its nullspace contains a subspace of
dimension equal to the number of columns minus \(D\).  Let \(P\) be
the orthogonal projection onto such a subspace.
Then \(\|H-Q\|_F^2\ge\|(H-Q)P\|_F^2=\|HP\|_F^2\).
In an eigenbasis of \(H^TH\), the latter is a weighted sum of its
eigenvalues with weights in \([0,1]\) summing to the rank of \(P\).
It is at least the sum of the smallest required eigenvalues,
which is \(T_D(H)\).  Keeping the largest \(D\) terms of a singular-value
decomposition attains this lower bound.  This is the
Eckart--Young approximation theorem~\cite{op:eckart}.
Finally, distance to any fixed nonempty subset is \(1\)-Lipschitz by
the triangle inequality; taking infima proves the last statement.
\end{proof}

The rank-\(D\) approximation just used is allowed to be any real matrix.
It need not be nonnegative, normalized, or realizable by a hidden process.

\begin{theorem}[Total variation bounds with marginal information]
\label{op:marginal-rank}
Let \(H,Q\) be normalized nonnegative joint matrices, with
\(\operatorname{rank}Q\le D\).  Write
\[
 \delta=\operatorname{TV}(H,Q),\quad
 r=\operatorname{TV}(H_X,Q_X),\quad
 c=\operatorname{TV}(H_Y,Q_Y),\quad T=T_D(H).
\]
Then
\[
 T\le\|H-Q\|_F^2\le
 \min\{2\delta^2,\delta(\delta+r+c)\},
\]
and therefore
\[
 \delta\ge
 \max\left\{\sqrt{T/2},
 \frac{\sqrt{(r+c)^2+4T}-(r+c)}2\right\}.
\]
When both marginals match, the second bound becomes \(\delta\ge\sqrt T\).
\end{theorem}
\begin{proof}
Put \(A=H-Q\).  Its positive and negative entry masses both equal
\(\delta\).  The sum of squares within each sign is at most the square
of its total mass, proving \(\|A\|_F^2\le2\delta^2\).
For a positive entry \(A_{ij}=x\), the negative mass in row \(i\) is at
least \(x-r\), since the positive part of any row sum is at most \(r\).
The negative mass in column \(j\) is at least \(x-c\).
These sets of negative entries are disjoint, because their only possible
intersection is the positive entry \((i,j)\).  Consequently
\(\delta\ge2x-r-c\).  Applying the same argument to \(-A\) gives
\(\|A\|_{\max}\le(\delta+r+c)/2\).
Thus \(\|A\|_F^2\le\|A\|_{\max}\|A\|_1
\le\delta(\delta+r+c)\).
Lemma~\ref{op:tail} supplies \(T\le\|A\|_F^2\);
solving the two quadratic inequalities proves the result.
\end{proof}

Known upper bounds on \(r+c\) may be substituted, because the second
lower bound decreases with that sum.  For zero row and column sums,
\(\|A\|_F\le\delta\) is sharp, as the signed \(2\)-by-\(2\) cycle
\(\left(\begin{smallmatrix}x&-x\\-x&x\end{smallmatrix}\right)\) shows.
That norm sharpness does not show that a rank-constrained probability
model attains the combined bound.

For \(L=2,\eta=1/2,\epsilon=1/10\), the exact noisy matrix is
\begin{equation}\label{op:noisy-four}
 \widetilde H_2=\frac1{2000}
 \begin{pmatrix}
 209&105&81&105\\
 81&105&209&105\\
 105&81&105&209\\
 105&209&105&81
 \end{pmatrix}.
\end{equation}
Its singular values are \(1/4,8/125,8/125,1/25\), and
\[
 T_1=153/15625,\qquad T_2=89/15625,\qquad T_3=1/625.
\]
Both marginals are uniform on four outcomes.  A normalized rank-one
nonnegative matrix with those same marginals is uniquely the matrix
whose entries are \(1/16\), and its total variation distance is exactly
\(21/125\).  The iid fair-bit process attains this restricted rank-one
competitor.  The number does not apply when competitor marginals may
vary; the bounds for two or three states have not been asserted attainable.

\subsubsection{From population rank to finite samples}

\begin{lemma}[Bounded-differences inequality]\label{op:bounded}
Let \(X_1,\ldots,X_n\) be independent, and suppose changing only \(X_i\)
changes \(f\) by at most \(c_i\).  Then
\[
 P(f-Ef\ge t)\le
 \exp\left(-\frac{2t^2}{\sum_i c_i^2}\right).
\]
\end{lemma}
\begin{proof}
The conditional Doob martingale of \(f\) has each martingale difference
in an interval of length at most \(c_i\): conditional expectations for
two possible values of \(X_i\) differ by at most \(c_i\), after coupling
the still independent remaining coordinates identically.
For a centered variable \(U\) in an interval of length \(c\),
\(\log E e^{\lambda U}\le\lambda^2c^2/8\).
Indeed the second derivative of its log moment-generating function is
the variance under exponential tilting, which is at most \(c^2/4\);
the function and its first derivative vanish at zero, so integration
twice gives the bound.  The variance bound follows from
\(E[(U-a)(b-U)]\ge0\), which gives
\(\operatorname{Var}(U)\le(EU-a)(b-EU)\le(b-a)^2/4\).
Apply the exponential bound conditionally to the martingale differences,
iterate, use Markov's inequality, and minimize over
\(\lambda=4t/\sum c_i^2\).  This is the standard bounded-differences
method~\cite{op:mcdiarmid}.
\end{proof}

\begin{theorem}[Empirical matrix radius and rank test]\label{op:finite-rank}
Fix a matrix shape with \(K\) cells, an observation map, and a sample size
\(n\).  Observe \(n\) independent records having the same categorical
joint law \(H\), and let \(\widehat H\) be their frequency matrix.
Dependence within a record is allowed.  For \(0<\alpha<1\), set
\[
 \varepsilon_\alpha(n,K)=
 \frac{\sqrt{1-1/K}+\sqrt{\log(1/\alpha)}}{\sqrt n}.
\]
Then \(P(\|\widehat H-H\|_F\le\varepsilon_\alpha)\ge1-\alpha\).
The test rejecting rank at most \(D\) when
\(d_D(\widehat H)>\varepsilon_\alpha\) has size at most \(\alpha\).
If \(d_D(H)>\varepsilon_\alpha+\varepsilon_\beta\), its power is at
least \(1-\beta\).
On the same confidence event, simultaneously for every \(D\),
\[
 \inf_{\substack{Q\ge0,\ \sum Q=1\\\operatorname{rank}Q\le D}}
 \operatorname{TV}(H,Q)\ge
 \frac{[d_D(\widehat H)-\varepsilon_\alpha]_+}{\sqrt2}.
\]
\end{theorem}
\begin{proof}
The multinomial second-moment identity gives
\[
 E\|\widehat H-H\|_F^2
 =\frac{1-\|H\|_F^2}{n}\le\frac{1-1/K}{n}.
\]
Jensen's inequality bounds the mean norm by the square root.
Replacing one record changes \(\widehat H\) by Frobenius norm at most
\(\sqrt2/n\).  The reverse triangle inequality gives the same
sensitivity for its error norm.  Lemma~\ref{op:bounded} therefore gives
an upper tail \(\exp(-nt^2)\); use
\(t=\sqrt{\log(1/\alpha)/n}\).
Under the rank null, \(d_D(H)=0\); Lemma~\ref{op:tail} proves the size
bound.  On the radius event with probability at least \(1-\beta\),
\(d_D(\widehat H)\ge d_D(H)-\varepsilon_\beta>\varepsilon_\alpha\),
which proves power.  No additional \(\alpha+\beta\) penalty is needed
for that power statement.  Finally,
\(d_D(H)\ge[d_D(\widehat H)-\varepsilon_\alpha]_+\), and
Theorem~\ref{op:marginal-rank}'s unrestricted bound applies to every
normalized rank-\(D\) competitor.
\end{proof}

At the single predeclared matrix and horizon, the confidence event
simultaneously controls all \(D\); repeated looks at changing horizons
are a different experiment.  Complete categorical outcomes, including
any failure outcome required by the experiment, must be represented.
A numerically computed singular tail requires an error allowance if
rounding can affect the threshold.

For \eqref{op:noisy-four}, \(K=16\) and \(\alpha=\beta=1/20\).
Let \(C=\sqrt{15/16}+\sqrt{\log20}\).
The condition is \(n>4C^2/T_D\).  Its first integer solutions are
\[
 \begin{array}{c|ccc}
 D&1&2&3\\ \hline
 n&2976&5116&18213\\
 \text{raw outputs }4n&11904&20464&72852
 \end{array}
\]
These are sufficient numbers of iid four-output records for this
concentration bound, not optimal testing budgets.  At \(n=18213\) the
single common power event suffices for all three rank rejections.
Experiments separately stopped at the three listed horizons do not
automatically constitute one test of size \(0.05\).

\subsection{Sequential evidence for complete record laws}
\label{op:sequential}

The null in this subsection is the complete iid fair-bit law,
\(P_0(Y_{1:n}=y)=2^{-n}\).
Fair one-bit marginals alone do not define this null.
Fix coherent candidate laws \(P_\theta\) before seeing data and
nonnegative mixing weights totaling one.

\begin{theorem}[Likelihood mixtures and an anytime threshold]
\label{op:mixture}
Each \(R_{\theta,n}=2^nP_\theta(Y_{1:n})\) is a nonnegative
\(P_0\)-martingale starting at one.  Their fixed mixture
\(\mathcal E_n=\int R_{\theta,n}\,dW(\theta)\) is also such a martingale,
and
\[
 P_0\{\exists n\ge0:\mathcal E_n\ge1/\alpha\}\le\alpha.
\]
\end{theorem}
\begin{proof}
Conditional on a prefix under \(P_0\), the next bit is fair.
Coherence gives
\[
 E_0[R_{\theta,n+1}\mid Y_{1:n}]
 =2^n\sum_{b=0}^1P_\theta(Y_{1:n},b)=R_{\theta,n}.
\]
Nonnegative integration permits conditional Tonelli, proving the mixture
claim.  Stop at the first threshold crossing or at a fixed \(N\).
The expected stopped value is one by finite optional stopping, so
\(P_0(\text{crossing by }N)\le\alpha\).
Let \(N\) increase.  The same argument holds for nonnegative
supermartingales and is Ville's inequality; see
Howard et al.~\cite{op:howard}.
\end{proof}

For the noisy parity family with fixed \(\eta=1/2,\epsilon=1/10\)
and unknown \(L\), take \(w_L=1/[L(L+1)]\).
When \(L\ge n\), every \(n\)-bit observed law is uniform.
Since \(\sum_{L=n}^\infty w_L=1/n\), the mixture has the finite expression
\begin{equation}\label{op:length-mixture}
 \mathcal E_n=\frac1n+
 \sum_{L=1}^{n-1}\frac{2^nP_L(Y_{1:n})}{L(L+1)},\qquad n\ge1,
 \quad \mathcal E_0=1.
\end{equation}
This identity does not guarantee inexpensive evaluation: a component
may require \(2^L\) hidden states.  Replacing the fixed mixture by a
post-data maximum does not preserve its martingale proof.

\begin{proposition}[A sufficient \(1950\)-output horizon]\label{op:1950}
At the fixed true alternative \(L=2,\eta=1/2,\epsilon=1/10\), the test
crossing \(20\) in \eqref{op:length-mixture} has power by \(1950\)
raw outputs at least \(0.9519745696463602\ldots\).
\end{proposition}
\begin{proof}
The \(650\) disjoint three-bit parity windows have independent
even-parity probability \(p=157/250\).  Let \(K\) be their even count and
let
\[
 S=(157/125)^K(93/125)^{650-K}
\]
be their likelihood ratio against fair parity.
For the full-record component \(R=2^{1950}P_2(Y_{1:1950})\),
strict positivity gives
\(E_2(S/R)=E_0S=1\).
The mixture contains \(R/6\).  Failure to cross by this horizon implies
in particular \(R<120\).
For any \(c>0\),
\[
 P_2(R<120)\le P_2(S<120c)+P_2(S/R>c)
              \le P_2(S<120c)+1/c .
\]
With \(c=40\), exact integer comparisons show
\(S\ge4800\) if and only if \(K\ge384\).
Therefore the miss probability is at most
\[
 \sum_{k=0}^{383}\binom{650}{k}
       (157/250)^k(93/250)^{650-k}+\frac1{40}
 =0.04802543035363984\ldots<\frac1{20}.
\]
All quantities in the last inequality are rational.  The calculation
does not require evaluating all nonconstant mixture components.
\end{proof}

This is a sufficient horizon for one specified alternative, not a
minimum, an expected stopping time, or a bound uniform in \(L\).
It is a different experiment from the fixed-horizon \(183\)-innovation
count test.

\begin{theorem}[Pointwise eventual detection]\label{op:eventual}
Fix a nondegenerate parity candidate with \(L<\infty\), \(\eta\ne0\),
\(0\le\epsilon<1/2\), and positive mixture weight \(w_L\).
If that candidate generates the data, its likelihood mixture eventually
crosses every fixed finite threshold almost surely.  More precisely,
\[
 \liminf_{n\to\infty}\frac{\log\mathcal E_n}{n}
 \ge \frac{d_{\rm Ber}(p_L\Vert1/2)}{L+1}>0,\qquad
 p_L=\frac{1+\eta(1-2\epsilon)^{L+1}}2 .
\]
\end{theorem}
\begin{proof}
Use \(m=\lfloor n/(L+1)\rfloor\) disjoint parity windows, and let \(S_m\)
be their likelihood ratio.  The iid strong law gives
\(\log S_m/n\to d_{\rm Ber}(p_L\Vert1/2)/(L+1)\).
For the full-record ratio \(R_{L,n}\), the same change-of-measure
identity as above gives \(E_L[S_m/R_{L,n}]=1\).
For every \(b>0\),
\[
 P_L(R_{L,n}<S_me^{-bn})\le e^{-bn}.
\]
The sum over \(n\) is finite, so the first Borel--Cantelli lemma
implies that this event occurs only finitely often.
Apply the assertion to a countable sequence of positive rational
\(b\) decreasing to zero and use
\(\mathcal E_n\ge w_LR_{L,n}\).
The fixed multiplicative weight has vanishing logarithm per \(n\).
\end{proof}

No uniform finite horizon follows: for \(L\ge n\) the full length-\(n\)
law equals the null exactly.  Alternatives approaching \(\eta=0\) or
\(\epsilon=1/2\) also approach the null on each fixed horizon.
For a concrete warning about the null, the process \(Y_t=B\) for one
fair bit \(B\) has fair marginals.  In the ideal-readout
\(\eta=1/2\) mixture, its \(L=1\) contribution is
\(\frac12(3/2)^{n-1}\), exceeding \(20\) at \(n=11\).
It is not an iid-null process; rejecting it does not violate the theorem.

\subsubsection{Conditional null envelopes and useful growth}

Suppose instead that the null supplies a predictable envelope
\[
 p_t=P(Y_t=1\mid\mathcal F_{t-1})
 \in[1/2-\delta_t,1/2+\delta_t],\quad 0\le\delta_t\le1/2 .
\]
For a predictable Bernoulli forecast \(q_t\), put \(q_t(1)=q_t\) and
\(q_t(0)=1-q_t\).  The factor
\[
 F_t=\frac{2q_t(Y_t)}
           {1+2\delta_t|2q_t-1|}
\]
has conditional null expectation at most one: its numerator has
expectation \(1+2(p_t-1/2)(2q_t-1)\), maximized exactly at the appropriate
endpoint of the interval.  Its products obey the crossing guarantee of
Theorem~\ref{op:mixture}.  This sharp normalization does not ensure
favorable evidence growth.

\begin{theorem}[Log-optimal binary factor for an interval null]
\label{op:clipped}
For \(0<\ell\le u<1\) and an alternative \(0<q<1\), let
\(p_*=\min\{u,\max\{\ell,q\}\}\).  Define
\[
 G_q(1)=q/p_*,\qquad G_q(0)=(1-q)/(1-p_*).
\]
It has expectation at most one for every Bernoulli parameter in
\([\ell,u]\), and maximizes \(E_q\log g(Y)\) over all nonnegative
factors \(g\) valid throughout that interval.
\end{theorem}
\begin{proof}
If \(q>u\), the expectation of \(G_q\) increases with the actual
Bernoulli parameter and equals one at \(u=p_*\).
If \(q<\ell\), it decreases and equals one at \(\ell=p_*\).
If \(q\in[\ell,u]\), it is identically one.
For every valid \(g\), Jensen's inequality gives
\[
 E_q\log\frac{g(Y)}{G_q(Y)}
 \le\log E_q\frac{g(Y)}{G_q(Y)}
 =\log E_{p_*}g(Y)\le0.
\]
Zero factors are handled by the extended value \(-\infty\) for the log.
\end{proof}

This is a binary specialization of the established growth-optimal
e-variable framework~\cite{op:larsson}.  Predictable choices of
\(\ell_t,u_t,q_t\) again yield a nonnegative null supermartingale.
The criterion is expected one-step log growth, not finite-horizon
Neyman--Pearson power.

For \(q=13/20\) and null interval \([2/5,3/5]\), the two factors are
\[
 (F(0),F(1))=(35/53,65/53),\qquad
 (G(0),G(1))=(7/8,13/12).
\]
Their expected log growth under iid Bernoulli\((13/20)\) is negative
and positive, respectively.  The exact sign certificates are
\[
 65^{13}35^7<53^{20},\qquad 13^{13}7^7>12^{13}8^7.
\]
The strong law makes the first product tend to zero and the second
diverge almost surely.  The first can still cross a threshold early.

A conditional-bias envelope is not a memorylessness null.  In the noisy
parity model, conditioning first on the full source history and then
averaging over the observed past yields
\[
 \left|P(\widetilde Y_t=1\mid\widetilde Y_{<t})-\frac12\right|
 \le\frac{|\eta|(1-2\epsilon)}2 .
\]
For the numerical alternative this is \(1/5\).
A null envelope at least that wide contains the dependent process, so
a test valid for every law in it must retain its null error guarantee
on that process.  The envelope itself needs justification for the
actual observation filtration; marginal calibration does not supply it.

\subsection{Visibility, environmental records, and access to fragments}
\label{op:records}

Consider an ordinary qubit and an environment with a fixed tensor
decomposition into \(N\) fragments.  A stipulated controlled interaction
maps the two qubit basis states to branch records
\[
 |0\rangle\longmapsto |0\rangle|e_0\rangle^{\otimes N},\qquad
 |1\rangle\longmapsto |1\rangle|e_1\rangle^{\otimes N},
 \quad \langle e_0|e_1\rangle=\gamma,\quad |\gamma|\le1.
\]
For a balanced input, tracing out the fragments leaves coherence of
magnitude \(|\gamma|^N/2\).  A phase setting with the convention
\(p_0(\phi)=\frac12[1+\operatorname{Re}(e^{i\phi}\gamma^N)]\)
has visibility \(|\gamma|^N\) and qubit purity
\((1+|\gamma|^{2N})/2\).
These formulas assume the stated state preparation, tensor partition,
controlled interaction, and Born-rule measurement.  They derive no
macroscopic amplification mechanism.

\begin{proposition}[A classical realization of the binary visibility law]
\label{op:visibility-alias}
For real \(0\le\gamma\le1\), all final binary laws
\[
 p_0(N,\phi)=\tfrac12(1+\gamma^N\cos\phi),
 \qquad N\in\{0,1,\ldots\},\ \phi\in\mathbb R,
\]
are reproduced by a two-state classical Markov chain and a
setting-dependent binary response.
\end{proposition}
\begin{proof}
Let \(Z_0=+1\).  Flip the sign independently at each step with
probability \((1-\gamma)/2\), so that the transition matrix is
\[
 K_\gamma=\frac12
 \begin{pmatrix}1+\gamma&1-\gamma\\1-\gamma&1+\gamma\end{pmatrix}.
\]
Then \(E Z_N=\gamma^N\), by the one-step conditional expectation
\(E(Z_{n+1}\mid Z_n)=\gamma Z_n\).
Given \(Z_N=z\) and the calibrated setting \(\phi\), report outcome zero
with probability \((1+z\cos\phi)/2\).  Averaging proves the formula.
At \(N=0\) the chain has made no transition, including when \(\gamma=0\).
\end{proof}

For complex \(\gamma=re^{i\beta}\), one may instead choose independent
classical phases \(\Theta_j=\beta\) or \(\beta+\pi\) with probabilities
\((1+r)/2,(1-r)/2\).  Since \(E e^{i\Theta_j}=\gamma\), the response
\([1+\cos(\phi+\sum_j\Theta_j)]/2\) has the required averaged law.
This construction uses accumulated phase; it is not the same assertion
of a homogeneous two-state realization with one fixed response for all
complex \(\gamma\).
Related random-unitary descriptions of quantum noise are established
\cite{op:crow}; their scope differs from this explicitly classical
simulation of a specified final binary family.  Neither construction
simulates arbitrary quantum preparations, joint measurements, Bell
experiments, or every environmental observable.

\begin{theorem}[A final channel does not identify fragment factorization]
\label{op:environment}
Fix the physical fragment partition and let \(0<\gamma<1\),
\(s=\sqrt{1-\gamma^2}\).  For \(N\ge2\), the product branches
\[
 E_0=|0\rangle^{\otimes N},\qquad
 E_1=(\gamma|0\rangle+s|1\rangle)^{\otimes N}
\]
and the branches \(F_i=U_NE_i\), where
\(U_N=\operatorname{CZ}_{12}\otimes I_{3:\!N}\), induce exactly the
same complete final qubit channel, including its action with a retained
ancilla.  Nevertheless \(F_1\) is not a product across the first two
fragments.
If the transformed old fragments are never accessed or recoupled, the
two descriptions also give the same complete histories for common
operations on the retained system, ancillas, and fresh fragments.
\end{theorem}
\begin{proof}
For any joint state \(\rho\) of a retained system and environment,
\[
 \operatorname{Tr}_E[(I\otimes U_N)\rho(I\otimes U_N^\dagger)]
 =\operatorname{Tr}_E\rho .
\]
This identity is valid with an arbitrary retained ancilla, proving the
complete-channel statement.  It also preserves the branch overlap
\(\langle E_0|E_1\rangle=\gamma^N\).
For the first two fragments, the coefficient matrix of \(F_1\) is
\[
 \begin{pmatrix}\gamma^2&\gamma s\\
                 \gamma s&-s^2\end{pmatrix},
 \qquad \det=-2\gamma^2s^2\ne0 .
\]
A pure two-fragment state is a product precisely when its coefficient
matrix has rank one; thus this branch is not a product.
Appending the remaining product factors does not remove that fact.
For a later history satisfying the stated access restriction, the
common environmental unitary commutes with every later retained-system
or fresh-fragment operation.  Move it through the product of the
corresponding instrument maps and remove it under the final trace.
This proves equality for each fixed outcome/action history; common
history-dependent randomized policies preserve it by conditioning.
\end{proof}

The history assertion requires absence of later coupling to the
transformed fragments.  Equality of final reduced channels by itself
does not prove equality under arbitrary later access.  Also, the
distinguished property here is tensor-product factorization of
conditional quantum fragment states.  It is not a claim that every
chosen classical fragment readout becomes statistically dependent.
For example, the common controlled phase leaves the branch's
\(Z\)-basis probability distribution unchanged.  If arbitrary entangling
changes of tensor factorization are declared to be a gauge freedom,
the target factorization property itself changes.
At \(\gamma=0\) or \(1\), or with fewer than two fragments, the stated
nonproduct witness is absent; common-unitary invariance of the reduced
channel still holds.

\begin{example}[A calibrated fragment measurement separates the pair]
\label{op:fragment-readout}
For \(\gamma=3/5,s=4/5,N=2\), the branch-one vectors are
\[
 E_1=(9,12,12,16)^T/25,\qquad
 F_1=(9,12,12,-16)^T/25.
\]
Both are normalized and have overlap \(9/25\) with \(|00\rangle\).
Their coefficient determinants are zero and \(-288/625\).
The first-fragment reduced off-diagonal entries are \(12/25\) and
\(-84/625\), so a local \(X\) measurement has
\[
 P_{E_1}(X=+)=49/50,\qquad P_{F_1}(X=+)=457/1250,
 \qquad \text{gap}=384/625.
\]
This requires anchored preparation of branch one, access to the specified
fragment, and a calibrated \(X\) readout.  If each reported binary law
is within total variation \(\varepsilon\) of its stated ideal law,
the gap is at least \(384/625-2\varepsilon\); strict positivity is a
sufficient separation condition.  This is pair separation, not
identification of an arbitrary environmental model.
\end{example}

\subsection{Positive hidden realizations and exact intervention contracts}
\label{op:controls}

A finite transition-emitting, or Mealy, model has nonnegative matrices
\(T_y\), indexed by a finite output alphabet, whose sum \(P=\sum_y T_y\)
is row-stochastic.  From an initial row distribution \(\pi\),
\[
 P(Y_{1:m}=y_1\cdots y_m)=\pi T_{y_1}\cdots T_{y_m}\mathbf1.
\]
A silent control is a row-stochastic matrix inserted between output
steps, producing no extra recorded outcome.  Linear reachability means
that the rows \(\{\pi T_w\}\) span the whole row space; linear
observability means that \(\{T_w\mathbf1\}\) spans the whole column
space, where \(w\) includes the empty word.  We use \emph{linearly minimal}
when both properties hold.  It is stronger than merely displaying no
obvious redundant state.

If \(S\) is invertible with \(S\mathbf1=\mathbf1\), put
\[
 T'_y=S^{-1}T_yS,\qquad \pi'=\pi S .
\]
Whenever these transformed quantities are valid probabilities, every
passive word law is unchanged by telescoping.  If each common fixed
control commutes with \(S\), all controlled word laws are unchanged as
well.  If instead controls are themselves allowed to vary, simultaneous
transformation \(K'_j=S^{-1}K_jS\) gives the same telescoping identity.
The distinction between exact fixed controls and uncertain controls
is substantive.  Similarity equivalence for Mealy models and its
distinction from other hidden-model classes are established
\cite{op:vanluyten}.

\begin{example}[A strictly positive minimal passive pair]
\label{op:mealy-pair}
Take
\[
 T_0=\frac1{20}\begin{pmatrix}6&2\\1&3\end{pmatrix},\quad
 T_1=\frac1{20}\begin{pmatrix}4&8\\7&9\end{pmatrix},\quad
 \pi=(4/9,5/9),\quad
 S=\frac1{10}\begin{pmatrix}9&1\\2&8\end{pmatrix}.
\]
The transformed matrices are
\[
 T'_0=\frac1{1400}\begin{pmatrix}449&151\\19&181\end{pmatrix},
 \qquad
 T'_1=\frac1{280}\begin{pmatrix}67&93\\125&115\end{pmatrix},
 \qquad \pi'=\pi S .
\]
Both representations are strictly positive and stationary:
\(\pi P=\pi\) and \(\pi'P'=\pi'\).
Their common two-output matrix is
\[
 J=\frac1{900}\begin{pmatrix}81&179\\179&461\end{pmatrix},
 \qquad \det J=53/8100.
\]
The factorization of \(J\) through a hidden boundary state implies
minimal dimension at least two; the displayed models attain it.
Their total-transition diagonal multisets are
\(\{1/2,3/5\}\) and \(\{14/25,27/50\}\), respectively.
They therefore are not related by a permutation of hidden-state labels.
\end{example}

\begin{proposition}[Exactly which fixed controls separate that pair]
\label{op:pair-control}
Let the same anchored control
\[
 K(a,b)=\begin{pmatrix}1-a&a\\b&1-b\end{pmatrix},
 \qquad 0\le a,b\le1,
\]
act in both models of Example~\ref{op:mealy-pair}.
All controlled output-history laws agree if and only if \(b=2a\).
The next-zero contrast after one such control is
\[
 \pi KT_0\mathbf1-\pi'KT'_0\mathbf1=(2a-b)/35.
\]
\end{proposition}
\begin{proof}
Direct multiplication gives the contrast and
\[
 SK-KS=\frac1{10}
 \begin{pmatrix}b-2a&2a-b\\b-2a&2a-b\end{pmatrix}.
\]
If \(b=2a\), the similarity telescopes through every interleaving
of outputs and controls.  If \(b\ne2a\), the displayed single-output
experiment already distinguishes the pair.  Common randomized or
adaptive control policies preserve the equality in the first case by
conditioning on their observed histories.
\end{proof}

The reset \(R_2=\left(\begin{smallmatrix}0&1\\0&1\end{smallmatrix}\right)\)
maximizes the absolute single-output contrast over this control square:
it gives \(2/35\), with zero probabilities \(1/5\) and \(1/7\).
This is an optimum for that contrast and that fixed pair.
It does not identify the entire model class.  For example,
\[
 U=\begin{pmatrix}9/10&1/10\\0&1\end{pmatrix}
 \quad\hbox{satisfies}\quad UR_2=R_2U=R_2 .
\]
The further transformed model has
\[
 \pi U=(2/5,3/5),\quad
 U^{-1}T_0U=\frac1{1800}\begin{pmatrix}531&229\\81&279\end{pmatrix},
 \quad
 U^{-1}T_1U=\frac1{1800}\begin{pmatrix}297&743\\567&873\end{pmatrix}.
\]
It is strictly positive, remains minimal, and has total-transition
diagonal multiset \(\{23/50,16/25\}\).
It is a further nonpermutation partner agreeing with the original on
all passive and exact-\(R_2\) controlled laws.

\begin{theorem}[One silent control does not identify the unrestricted
positive minimal class]\label{op:single-control}
Let \(n\ge2\), let \((\pi,\{T_y\})\) be a strictly positive linearly
minimal finite-alphabet realization, and assume its unanchored
preparation and output matrices may vary in a neighborhood of their
current values within the model class.  A single exact common
row-stochastic silent control \(K\) leaves a continuum of distinct
positive realizations with the same complete controlled laws.
Stationarity, if imposed on \(\pi\), is preserved.
\end{theorem}
\begin{proof}
For \(K\ne I\), choose \(S_t=I+t(K-I)\).
It satisfies \(S_t\mathbf1=\mathbf1\), commutes with \(K\), and is
invertible for sufficiently small \(t\).
At \(t=0\), all entries of \(T_y\) and \(\pi\) are positive.
Continuity and finiteness preserve positivity for a sufficiently small
interval of \(t\).  Normalization follows from \(S_t\mathbf1=\mathbf1\).
The transformed total transition matrix and preparation satisfy
\(\pi'P'=\pi P S_t=\pi S_t=\pi'\) when the original is stationary.
The telescoping identity proves equality for all controlled words.
If \(K=I\), use any nonzero matrix \(A\) with \(A\mathbf1=0\) and
\(S_t=I+tA\), again sufficiently near the identity.
To see that distinct similarities give distinct numerical realizations,
suppose a similarity \(V\) fixes both \(\pi\) and every \(T_y\).
Then \(\pi T_wV=\pi T_w\) for every word.  Reachability forces \(V=I\).
The ratio of two similarities yielding identical realizations would
be such a \(V\), hence the similarities would be equal.
A continuum of distinct realizations cannot be removed by the finite
set of state permutations.
\end{proof}

The theorem concerns a local open positive class, an unanchored
preparation, and silent stochastic controls.  Structural zeros, fixed
preparations, recorded control outcomes, or narrower parametrizations
can change the conclusion.  It is not a claim about all hidden-model
classes.  Adding invisible hidden labels is an easier ambiguity:
one may append arbitrary hidden transitions whose pushed visible rows
stay fixed, as in Theorem~\ref{op:lumpability}.  Such redundant lifts
are not counterexamples to minimal-realization theorems and do not
establish distinct physical ontologies.

\begin{theorem}[A sharp two-control criterion in two states]
\label{op:two-controls}
Consider stationary strictly positive binary Mealy models with two
states, passive linear dimension two, and an unanchored initial
distribution.  Suppose two labelled controls are known exactly and
have the same matrices in every admissible model:
\[
 K_j=\begin{pmatrix}1-a_j&a_j\\b_j&1-b_j\end{pmatrix},\quad j=1,2 .
\]
Complete controlled population laws identify every numerical output
matrix and the labelled preparation in this class if and only if
\[
 \Delta=a_1b_2-a_2b_1\ne0 .
\]
When \(\Delta\ne0\), experiments with at most one control followed by
at most two outputs suffice for the reconstruction.
\end{theorem}
\begin{proof}
Write \(\pi=(r,1-r)\), \(e=T_0\mathbf1=(e_1,e_2)^T\),
\(d=e_1-e_2\), \(k=rd\), \(q=P(0)\), and
\(q_j=P(K_j\text{ then }0)\).
If \(d=0\), both \(T_0\mathbf1=q\mathbf1\) and
\(T_1\mathbf1=(1-q)\mathbf1\).  Repeated multiplication makes every
word iid Bernoulli, contradicting passive dimension two.  Thus \(d\ne0\).
The first-state mass after \(K_j\) is
\(r+b_j-(a_j+b_j)r\), hence
\[
 \begin{pmatrix}q_1-q\\q_2-q\end{pmatrix}
 =
 B\begin{pmatrix}d\\k\end{pmatrix},\qquad
 B=\begin{pmatrix}
 b_1&-(a_1+b_1)\\ b_2&-(a_2+b_2)
 \end{pmatrix},\quad \det B=\Delta .
\]
If \(\Delta\ne0\), solve for \(d,k\), then recover
\[
 r=k/d,\qquad e_2=q-k,\qquad e_1=e_2+d .
\]
At least one \(q_j-q\) is nonzero, because otherwise \(B\) would map
the nonzero vector \((d,k)^T\) to zero.  For such \(j\), define
\[
 A=\begin{pmatrix}\pi\\\pi K_j\end{pmatrix},\quad
 F=\begin{pmatrix}\mathbf1&e\end{pmatrix},\quad
 C_y=\begin{pmatrix}
 P(y)&P(y0)\\
 P(K_j\text{ then }y)&P(K_j\text{ then }y0)
 \end{pmatrix}.
\]
Then \(C_y=AT_yF\),
\(\det A=-(q_j-q)/d\), and \(\det F=-d\).
Both are invertible, so \(T_y=A^{-1}C_yF^{-1}\).
All quantities are obtained from the stated short experiments.

Conversely, direct multiplication gives
\[
 [K_1,K_2]=\Delta\begin{pmatrix}1&-1\\1&-1\end{pmatrix}.
\]
If \(\Delta=0\) and a control \(K\ne I\) exists, the nonconstant
similarities \(I+t(K-I)\) commute with both controls.
If both controls are identity, use any nontrivial near-identity
row-normalized similarity.  The positivity, stationarity, and
minimality argument of Theorem~\ref{op:single-control} provides a
continuum of admissible distinct realizations with all controlled
laws equal.  Identification is impossible.
\end{proof}

For a positive non-reset example,
\[
 K_1=\begin{pmatrix}3/4&1/4\\1/2&1/2\end{pmatrix},\qquad
 K_2=\begin{pmatrix}1/2&1/2\\1/4&3/4\end{pmatrix},\qquad
 \Delta=-3/16.
\]
For Example~\ref{op:mealy-pair},
\[
 q=13/45,\quad q_1=29/90,\quad q_2=49/180,\quad
 d=1/5,\quad k=4/45,
\]
and the short-word matrices for \(j=1\) are
\[
 C_0=\begin{pmatrix}13/45&9/100\\29/90&21/200\end{pmatrix},
 \qquad
 C_1=\begin{pmatrix}32/45&179/900\\61/90&337/1800\end{pmatrix}.
\]
The formulas recover the displayed \(\pi,T_0,T_1\) exactly.
Small \(|\Delta|\), small \(|d|\), or a small preparation contrast can
make the inverse numerically unstable; this population theorem
contains no finite-sample stability bound.

\begin{proposition}[Pure preparations and a finite-state positive control]
\label{op:anchored-preparations}
For a two-state model with two exactly anchored labelled pure
preparations, put
\(e_i=P_i(0)\), \(p_i(y)=P_i(y)\), \(z_i(y)=P_i(y0)\).
If \(d=e_1-e_2\ne0\), then
\[
 (T_y)_{i1}=\frac{z_i(y)-e_2p_i(y)}d,\qquad
 (T_y)_{i2}=\frac{e_1p_i(y)-z_i(y)}d.
\]
More generally, an observable \(n\)-state realization is identified by
an exact labelled cyclic permutation control and an exact reset to
one labelled state, using experiments with at most \(n\) output symbols.
\end{proposition}
\begin{proof}
For two states, the equations
\((T_y)_{i1}+(T_y)_{i2}=p_i(y)\) and
\(e_1(T_y)_{i1}+e_2(T_y)_{i2}=z_i(y)\)
give the formulas.
For \(n\) states, a reset followed by powers of the cycle prepares
every pure labelled state.  Let
\(V_k=\operatorname{span}\{T_v\mathbf1:|v|\le k\}\).
The initial space \(V_0\) has dimension one.
If \(V_k=V_{k+1}\), each \(T_y\) maps \(V_k\) into itself, so all later
spaces are identical.  Observability therefore implies that full
dimension is reached by \(k=n-1\).
Choose \(n\) spanning suffix columns to form an invertible matrix \(F\).
Their probabilities under the pure preparations give every entry of
\(F\).  Probabilities of \(yv\) under the same preparations give
\(T_yF\), so recover \(T_y\).  Probabilities of those suffixes from the
original preparation give \(\pi F\), so recover \(\pi\).
\end{proof}

The cycle and reset contain structural zeros and permit repeated uses.
The bound concerns the number of output symbols per setting; it does
not count control uses, distinct settings, repeated trials, duration,
or calibration cost.  Together with the one-control obstruction, two
types are minimal only for the stated universal ideal control-design
question over positive minimal models.  Restricted model classes may
require fewer.

\subsection{Calibration intervals and statistical separation}
\label{op:calibration}

Exact controls and approximately known controls define different
inverse problems.  This can be shown quantitatively without assigning
a prior to the uncertain control.

\begin{theorem}[An explicit calibration-induced identification interval]
\label{op:calibration-interval}
Use the positive model of Example~\ref{op:mealy-pair} and the two
controls displayed after Theorem~\ref{op:two-controls}.
Keep \(K_1\) fixed exactly, but allow \(K_2\) to lie in its full
maximum-row total-variation ball of radius \(\varepsilon\).
Keep the preparation unanchored.  Among strictly positive stationary
models with the same complete controlled output laws, the first
preparation coordinate includes the interval
\[
 \left[\frac49,\frac49+\min\left\{\frac1{15},
                                      \frac{8\varepsilon}{9}\right\}\right].
\]
This is a contained interval, not the complete identified set.
\end{theorem}
\begin{proof}
For \(0\le t\le1/10\), use
\[
 S_t=\begin{pmatrix}1-t&t\\2t&1-2t\end{pmatrix},
 \qquad \det S_t=1-3t\ge7/10.
\]
Direct multiplication gives
\[
 S_t^{-1}K_1S_t=K_1,\qquad
 S_t^{-1}K_2S_t-K_2
 =\begin{pmatrix}3t/4&-3t/4\\3t/4&-3t/4\end{pmatrix}.
\]
Thus the maximum-row total variation error is \(3t/4\).
The transformed preparation is
\(\pi S_t=((4+6t)/9,(5-6t)/9)\).
For completeness, the eight numerators of the two transformed output
matrices, over the positive common denominator \(20(1-3t)\), are
\[
 \begin{array}{cc}
 6-15t-t^2&2-5t+t^2\\
 1-8t-t^2&3-12t+t^2\\[2pt]
 4-3t-35t^2&8-37t+35t^2\\
 7-4t-35t^2&9-36t+35t^2
 \end{array}.
\]
On the stated interval they are bounded below, in the same order, by
\[
 449/100,\ 3/2,\ 19/100,\ 9/5,\
 67/20,\ 43/10,\ 25/4,\ 27/5,
\]
all strictly positive.  Normalization and stationarity are preserved by
similarity.  Simultaneously transforming \(T_y,\pi,K_2\), while \(K_1\)
stays fixed, makes every controlled word telescope to its original
probability.  The admissible segment is
\(0\le t\le\min\{1/10,4\varepsilon/3\}\).
Its first preparation coordinate is \(4/9+2t/3\), giving the interval.
\end{proof}

For example, \(\varepsilon=1/100\) gives
\([4/9,34/75]\), of width \(2/225\).
This is uncertainty in a chosen latent coordinate, not uncertainty in
the common predictions.  An independently fixed preparation or a
model for extra calibration outcomes can break the equivalence.

\begin{proposition}[Finite positive menus with open calibration classes]
\label{op:open-calibration}
Let \(n\ge2\), with finitely many strictly positive output matrices,
a positive unanchored preparation, and finitely many strictly positive
silent controls.  Suppose their admissible classes contain local
neighborhoods of the output/preparation parameters and relative-open
neighborhoods of each control within row-stochastic matrices.
If the realization is linearly minimal, there is a local continuum of
distinct admissible simultaneous-similarity realizations with the same
complete controlled laws.
\end{proposition}
\begin{proof}
Choose a nonzero matrix \(A\) with \(A\mathbf1=0\), and
\(S_t=I+tA\) sufficiently near \(I\).
Transform every output matrix and every control by \(S_t^{-1}(\cdot)S_t\),
and the preparation by \(\pi S_t\).
The sum of the transformed output matrices and each transformed control
are row-stochastic, and the transformed preparation sums to one, because
\(S_t\mathbf1=\mathbf1\).
Strict positivity and the finitely many open-neighborhood conditions
persist for sufficiently small \(t\).
All probabilities telescope.  The trivial stabilizer argument in
Theorem~\ref{op:single-control} proves distinctness.
\end{proof}

The open-neighborhood premise excludes exact structural zeros and other
constraints that obstruct the path.  This qualitative result provides
no universal numerical interval width.

\subsubsection{The reset experiment with bounded nuisance}

For the fixed pair in Example~\ref{op:mealy-pair}, suppose the intended
reset to state two may place a candidate-specific mass \(u\in[0,\delta]\)
on state one.  The next-zero laws range over
\[
 p_A=(1+u)/5,\qquad p_B=(1+2u)/7.
\]
For \(0\le\delta<1/5\), the two probability intervals have strictly
positive separation
\[
 g_\delta=\frac15-\frac{1+2\delta}7
          =\frac{2-10\delta}{35}.
\]
At \(\delta=1/5\), take the exact \(R_2\) in the first model and
\[
 S^{-1}R_2S=
 \begin{pmatrix}1/5&4/5\\1/5&4/5\end{pmatrix}
\]
in the second.  The complete controlled laws, not only the next output,
then agree by simultaneous similarity.  The same witness is admissible
for every larger \(\delta\).  Thus this fixed-pair nuisance class has
a sharp \(20\%\) threshold for uniform separation by arbitrary records.
With an additional total-variation readout allowance \(\epsilon\) on
each binary law, \(g_\delta-2\epsilon>0\) is a sufficient one-output
separation condition.  The full-history sharpness claim above concerns
the explicitly specified reset nuisance, not every possible
history-dependent readout-error class.

With perfect resets, perfect readout, and independent repeated reset
trials, the simple hypotheses are Bernoulli\((1/7)\) versus
Bernoulli\((1/5)\) for the zero indicator.
Lemma~\ref{op:np} yields optimal size-\(1/20\) randomized powers
\[
 \begin{array}{c|cc}
 n&467&468\\ \hline
 \text{power}&0.9498952948087144\ldots&
               0.9505519541817742\ldots
 \end{array}
\]
At \(n=468\), reject for at least \(81\) zeros and randomize at \(80\)
with probability \(0.998625929855047\ldots\), defined exactly by
\[
 \gamma=\frac{1/20-
   \sum_{k=81}^{468}\binom{468}{k}(1/7)^k(6/7)^{468-k}}
   {\binom{468}{80}(1/7)^{80}(6/7)^{388}}.
\]
Discarding an extra iid trial cannot reduce attainable full-record
power, so the \(467/468\) boundary proves the minimum for this specified
randomized fixed-horizon test class.  It does not include calibration
trials or apply to the two-control reconstruction theorem.

For a different, robust experiment suppose the conditional zero
probability after every observed past is at most \(p_-\) under one
class and at least \(p_+\) under the other, with \(g=p_+-p_->0\).
Threshold the mean at \((p_-+p_+)/2\).
The conditional exponential bound used in
Lemma~\ref{op:bounded}, applied to the bounded martingale differences
of the zero indicators, gives each error at most
\(\exp(-ng^2/2)\).
It suffices that \(n\ge\lceil2\log(20)/g^2\rceil\).
For \(g=2/35\), this is \(1835\) trials.
Independence is unnecessary for this conditional-mean argument, but
the interval must hold under the actual filtration at every trial.
An unconditional mean gap does not supply it.

\subsection{Sampled generators and the information supplied by a clock}
\label{op:time}

A finite continuous-time Markov generator \(Q\) has nonnegative
off-diagonal entries and zero row sums, and produces transition
matrices \(e^{tQ}\) for \(t\ge0\).
Recovering \(Q\) from sampled transitions is the classical embedding
and matrix-exponential inverse problem
\cite{op:baake,op:yue}.  It presupposes a state basis and time coordinate.

\begin{theorem}[Two positive generators with all integer samples equal]
\label{op:generator-alias}
Let
\[
 C=\begin{pmatrix}0&1&0\\0&0&1\\1&0&0\end{pmatrix},\quad
 \Pi=\tfrac13\mathbf1\mathbf1^T,\quad R=I-\Pi,\quad A=C-C^2,
\]
and define
\[
 Q_0=-12R,\qquad Q_1=-12R+\frac{2\pi}{\sqrt3}A .
\]
Both are irreducible generators with uniform stationary distribution
and strictly positive off-diagonal rates.
They obey \(e^{nQ_0}=e^{nQ_1}\) for every integer \(n\ge0\), but are
neither similar nor related by a positive scalar time rescaling,
even when a similarity is also allowed.
One is reversible and the other has nonzero stationary circulation.
\end{theorem}
\begin{proof}
The off-diagonal rates are \(4\) for \(Q_0\) and
\(4\pm2\pi/\sqrt3>0\) for \(Q_1\).
The inequality follows, for example, from
\(\pi<22/7<2\sqrt3\).
Rows and columns sum to zero, so the uniform distribution is
stationary; strict positive rates imply irreducibility and uniqueness.
Direct algebra gives
\[
 A\Pi=\Pi A=0,\quad A^2=-3R,\quad AR=RA=A.
\]
Summing the exponential series on the invariant subspaces yields
\[
 e^{tQ_0}=\Pi+e^{-12t}R,\qquad
 e^{tQ_1}=\Pi+e^{-12t}
 \left[\cos(2\pi t)R+\frac{\sin(2\pi t)}{\sqrt3}A\right].
\]
At integer \(t\) the kernels agree.
The spectra are \(0,-12,-12\) and \(0,-12\pm2\pi i\), respectively.
Similarity preserves spectra, and a positive scalar cannot turn real
nonzero eigenvalues into nonreal ones.
The symmetric generator \(Q_0\) is reversible for the uniform law.
For \(Q_1\), the oriented stationary flux difference along a cycle edge
is \(\frac13[(4+2\pi/\sqrt3)-(4-2\pi/\sqrt3)]
=4\pi/(3\sqrt3)>0\).
\end{proof}

\begin{proposition}[Exactly the integer-grid intervention menu]
\label{op:grid-policy}
Use either generator in Theorem~\ref{op:generator-alias}, a common
initial law, and the same instantaneous labelled boundary controls
and instruments.  Suppose every control and every readout, including
silent state maps, occurs on the integer time grid, the between-boundary
generator stays fixed, and the policy depends only on the common
observed history and private randomness.
Then the complete finite controlled observation-history laws agree.
\end{proposition}
\begin{proof}
A fixed outcome-history probability is a row product of the initial
law, transition matrices for nonnegative integer durations, and the
corresponding boundary instrument matrices, followed by \(\mathbf1\).
Every factor is common to the two descriptions, so the probability
is common.  For an adaptive randomized policy, condition successively
on each history; its next action, scheduling choice, and stop decision
have the same conditional law.  Induction gives the result.
\end{proof}

Integer-grid readouts alone are insufficient if silent controls are
allowed at off-grid times.  Persistent drives that change the generator,
continuous monitoring, and observations of jump or waiting times also
lie outside this menu and may supply distinguishing information.

An anchored preparation of state one and a readout distinguishing
states two and three at time \(t=1/4\) give the contrast
\[
 (e^{Q_1/4})_{12}-(e^{Q_1/4})_{13}
 =\frac{2e^{-3}}{\sqrt3}
 =0.0574891546486971\ldots ,
\]
whereas the contrast is zero for \(Q_0\).
An unconditioned stationary uniform preparation has identical uniform
marginals under both models and does not provide this witness.
The contrast assumes a calibrated quarter-time setting; it is not
claimed optimal and carries no trial-count guarantee.

\begin{theorem}[Exact identification from two incommensurate times]
\label{op:two-times}
Let \(Q,B\) be finite real matrices represented in the same state basis.
Let \(t_1,t_2>0\) have irrational ratio.
If
\[
 e^{t_1Q}=e^{t_1B},\qquad e^{t_2Q}=e^{t_2B},
\]
then \(Q=B\).
\end{theorem}
\begin{proof}
Denote the two common exponential matrices by \(M_1,M_2\).
They are invertible and commute, since they are both functions of \(Q\).
For every pair of integers \(m,n\),
\[
 e^{(mt_1+nt_2)Q}=M_1^mM_2^n=e^{(mt_1+nt_2)B}.
\]
The subgroup \(\{mt_1+nt_2:m,n\in\mathbb Z\}\) is dense in
\(\mathbb R\).  One proof is that irrational rotation modulo \(t_2\)
has arbitrarily small positive return differences by the pigeonhole
principle; integer multiples of those differences approximate any
real number, and all such multiples remain in the subgroup.
Continuity gives equality of the exponentials for every real time.
Differentiating at zero yields \(Q=B\).
\end{proof}

The proof does not assume \(Q\) and \(B\) commute with one another.
Negative powers are algebraic inverses, not proposed physical operations
at negative time.  The result requires full exact transition matrices,
a common basis, homogeneous dynamics, and exact times.
Finite-precision measurements cannot establish exact irrationality of
a time ratio, and this uniqueness theorem is not a stability theorem.

\begin{theorem}[No uniform inverse stability over unbounded rates]
\label{op:rate-instability}
Fix a nonzero real \(c\) and any finite collection of strictly positive
observation times.  With the matrices \(R,A\) above, put
\[
 Q_a=-aR,\qquad B_a=-aR+cA,\qquad a>3|c|.
\]
These are irreducible Markov generators.  Their difference is the fixed
nonzero matrix \(cA\), while
\(\|e^{tQ_a}-e^{tB_a}\|\to0\) as \(a\to\infty\) at every specified
time, for any matrix norm.
Thus no uniform inverse error modulus holds over the unrestricted-rate
class.  In particular, even when the menu is injective, as for two
incommensurate times, its inverse is not uniformly continuous.
\end{theorem}
\begin{proof}
Off-diagonal entries are \(a/3\) or \(a/3\pm c\), all positive.
The same invariant-subspace calculation gives
\[
 e^{tB_a}-e^{tQ_a}
 =e^{-at}\left[
   (\cos(\sqrt3ct)-1)R+
   \frac{\sin(\sqrt3ct)}{\sqrt3}A\right].
\]
For fixed positive \(t\), the bracket is bounded independently of \(a\)
and the prefactor tends to zero.  A finite maximum over observation
times tends to zero as well, while
\(\|B_a-Q_a\|=|c|\|A\|>0\).
\end{proof}

On a specified compact generator class, the exact two-time map of
Theorem~\ref{op:two-times} is continuous and injective, so its inverse
on its image is uniformly continuous by compactness.
That qualitative exception provides no quantitative modulus, and no
joint finite-data result for uncertain timing and uncertain transition
matrices is obtained here.

\subsection{Scope of the mathematical conclusions}
\label{op:scope}

The positive results identify precisely what their declared experiments
determine: autonomous coarse dynamics under a fiber condition; a
classical predictive-state lower bound from a joint-law matrix; valid
finite-sample or sequential tests under explicit sampling assumptions;
finite-state reconstruction under exact anchored controls; and a
generator under full exact observations at two incommensurate times.
The counterexamples identify specific missing information: short-block
agreement, a final channel without fragment access, a single silent
control, open calibration uncertainty, or sampling on a coarse time
grid can leave distinct descriptions observationally equivalent.

All these conclusions are conditional on conventional mathematical
primitives and the stated observation maps.  They do not select a
physical tensor partition, demonstrate that a detector is independent,
derive a probability rule from a prephysical proposal, or provide an
emergent physical clock.  A physical correspondence must identify which
preparations, interventions, outcomes, reset procedures, and timings
realize the premises and must confront the appropriate competing
observation laws.
